% ============================================================================
% Section C: Distinguish Acquisition and Utilization (Deliverable C)
% ============================================================================
\section{Acquisition vs.\ Utilization Analysis}
\label{sec:C}

Life-cycle activities divide into two coupled domains:
\begin{itemize}
    \item \textbf{Acquisition-Oriented Activities (Phases 1--5):} Conceiving, permitting (CTE/CTO), building, and commissioning UA-EWRRF. Although CapEx represents only $\approx 25\%$ of total 15-year cost, acquisition decisions lock in over \textbf{70--80\%} of eventual life-cycle cost, safety risk, and maintainability.
    \item \textbf{Utilization-Oriented Activities (Phases 6--8):} Operating collection/recovery lines, maintaining availability, upgrading modules, and retiring the site cleanly. These phases represent over 72\% of total expenditure (power, wages, spares, TSDF gate fees, and closure).
\end{itemize}

Table~\ref{tab:annexC} examines three major acquisition-stage decisions whose impacts surface during utilization.

\begingroup
\footnotesize
\renewcommand{\arraystretch}{1.22}
\begin{xltabular}{\textwidth}{L{2.7cm} L{2.15cm} L{2.65cm} L{2.55cm} Y}
\caption{Acquisition-stage decisions significantly influencing utilization in UA-EWRRF (Annexure C format).}
\label{tab:annexC}\\
\toprule
\hdrrow \tblhead{Acquisition-stage decision} & \tblhead{Decision-maker} & \tblhead{Stakeholders experiencing consequences} & \tblhead{When consequences become visible} & \tblhead{Difficulty of changing later} \\
\midrule
\endfirsthead
\multicolumn{5}{l}{\footnotesize\color{SlateMuted}\textit{Table~\ref{tab:annexC} continued from previous page.}}\\
\toprule
\hdrrow \tblhead{Acquisition-stage decision} & \tblhead{Decision-maker} & \tblhead{Stakeholders experiencing consequences} & \tblhead{When consequences become visible} & \tblhead{Difficulty of changing later} \\
\midrule
\endhead
\bottomrule
\endfoot

\textbf{C1. Shred-first vs.\ pre-depollution hybrid sorting line} & Promoter Board \& EPC (Ph.\ 2--3) & Operators (fire/dust), buyers (impure lots), TSDF (toxic residue). & Years 1--2 via battery fires, alloy contamination, and buyer penalties. & \textbf{Very High:} Requires dismantling primary shredders and halting revenue for 3--5 months. \\
\midrule
\textbf{C2. Proprietary imported machines vs.\ standardized modular skids} & Procurement \& EPC Vendor (Ph.\ 3--4) & Maint.\ team (spare delays), promoter (downtime), OEM EPR clients. & Years 2--7 when imported blades/sensors take 6--10 weeks to replace. & \textbf{High:} Footprints, PLC controls, and power drops are locked to one vendor. \\
\midrule
\textbf{C3. Grouted concrete mounts vs.\ bolted skids on epoxy containment} & Civil EPC \& Cost Controller (Ph.\ 3--5) & Upgrade teams, MPCB, MIDC, decommissioning contractors. & Year 6 upgrades and Year 15 heavy-metal soil excavation audits. & \textbf{Extreme:} Re-pouring slabs beneath heavy machinery requires near-total demolition. \\

\end{xltabular}
\endgroup

\subsection{Stakeholder Alignment Conclusion}
Acquisition and utilization stakeholders are \textbf{only partially aligned}. Both want initial CTO approval and 10 t/day capacity, yet EPC contractors and CapEx-focused sponsors benefit from minimizing construction cost and schedule---favoring grouted mounts or monolithic machinery. Meanwhile, operations, maintenance, EHS staff, and retirement teams absorb the resulting 15-year downtime, dust exposure, retrofit blocks, and remediation liabilities. Resolving this gap requires LCC-based vendor bidding, 3-year OEM availability warranties, and early operator sign-off at PDR/CDR.
